\section{Learned closures}
\label{sec:learned}

Machine learning changes how a closure is represented and calibrated. It
does not remove the closure problem, and every structural question of the
previous sections applies unchanged. The literature chapter in
\texttt{shared/references/} treats the surrounding methods (PINNs, neural
operators, world models) at length; this section is restricted to closures
and to the design decisions that determine whether a learned closure can be
trusted. Reviews from the turbulence side are in
\cite{DuraisamyIaccarinoXiao19}.

\subsection{What is learned}

The first decision is the mathematical object. The options are the rows of
the table in Section~\ref{sec:kinds}, and the choice fixes what the model
can and cannot represent.
\begin{center}
\renewcommand{\arraystretch}{1.3}
\begin{tabularx}{\textwidth}{@{}P{3.6cm}P{5.2cm}Y@{}}
\toprule
Learned object & Form & Examples \\
\midrule
Coefficients of a classical closure
  & $\nu_t=g_\theta(\text{invariants})$
  & tensor-basis networks for the Reynolds stress \cite{Ling16} \\
The unclosed moment
  & $M_{N+1}=g_\theta(M_0,\dots,M_N)$
  & direct regression; no guarantee of hyperbolicity \\
Its gradient
  & $\partial_xM_{N+1}=\sum_ja_{j,\theta}\,\partial_xM_j$
  & hyperbolic moment closures \cite{Huang21,Christlieb25} \\
A nonlocal operator
  & $q=\mathcal W_\theta[T]$
  & learned Landau-fluid closures \cite{Ma20} \\
The choice of moments and the closure
  & encoder plus closed dynamics
  & learned generalized moments \cite{HanMaMaE19} \\
A memory kernel or auxiliary state
  & $K_\theta(s)$, or $\dot z=g_\theta(z,x)$
  & data-driven Mori--Zwanzig \cite{Lin21} \\
A collision operator
  & $C_\theta[f]$ in a conservative form
  & learned from molecular dynamics \cite{Zhao25} \\
A discretization-level correction
  & correction to a coarse solver step
  & learned corrections to coarse solvers \cite{Kochkov21} \\
\bottomrule
\end{tabularx}
\end{center}
Whole-solution surrogates (a network that maps initial data to the solution)
are not closures. They discard the reduced equation instead of completing
it, and in this project they are run as black-box controls.

\subsection{How structure enters}

A constraint can be imposed softly, as a penalty in the loss, or hard, by
construction. A soft constraint holds approximately on the training
distribution. A hard constraint holds everywhere and is the only kind that
supports a claim about behaviour in a long rollout.
\begin{itemize}
  \item \textbf{Conservation:} learn a flux, not a source, and apply a
    conservative difference to it.
  \item \textbf{Symmetry:} build inputs from invariants and outputs in a
    tensor basis (Galilean and rotational invariance).
  \item \textbf{Hyperbolicity:} parameterize the closure so that the flux
    Jacobian has real eigenvalues, for instance through a symmetrizer or
    through the eigenvalues themselves.
  \item \textbf{Dissipation:} parameterize a non-negative symbol or a
    positive semidefinite coefficient matrix.
  \item \textbf{Limits:} learn a correction to a classical closure that
    vanishes in the regime where the classical one is trusted,
    \begin{equation}
      q=q_{\mathrm{classical}}+\delta q_\theta ,
      \qquad \delta q_\theta\to0\ \text{as}\ k\lambda\to0 .
    \end{equation}
\end{itemize}
The last item is engineering principle 10 of this project: the classical
reduced operator stays explicit and the learned part supplies only an
identified correction. For the linear heat flux the results of
Section~\ref{sec:repo} make it concrete. In one velocity dimension the
correction satisfies $0\le\kF-\chi\le\kF Kk^2/(1+Kk^2)$, so its amplitude
and its long-wavelength scaling are known in advance. For a thermal drive
the exact symbol is $\int(m+k^2)^{-1}\dd\rho(m)$ with $\rho\ge0$, so a
finite sum $\sum_jw_j/(m_j+k^2)$ with non-negative weights has the right
analytic type, and two linear conditions on the weights fix the
conductivity and the leading deficit.

\subsection{How it is trained}

\begin{itemize}
  \item \textbf{A priori (offline):} regress the closure term on reference
    data. Cheap, and it optimizes the wrong quantity: the defect on
    reference states, not the behaviour of the closed model.
  \item \textbf{A posteriori (solver in the loop):} differentiate through
    the closed solver and penalize the error of the rolled-out solution
    \cite{ListChenThuerey22,Kochkov21}. It optimizes the right quantity,
    needs a differentiable solver, and couples the closure to one
    discretization.
  \item \textbf{Statistical targets:} fit long-time statistics when
    trajectories diverge chaotically and pointwise rollout error is
    meaningless.
\end{itemize}

\subsection{How it fails}

\begin{itemize}
  \item \textbf{The a priori to a posteriori gap.} A small defect on
    reference data, an unstable or biased closed model
    (Section~\ref{sec:turb}).
  \item \textbf{Distribution shift under rollout.} The closed model drifts
    to states outside the training set, where the closure is unconstrained.
  \item \textbf{Extrapolation in the regime parameter.} A closure trained
    at one collisionality, coupling strength or resolution is used at
    another. Figure~\ref{fig:between} and the electron results of
    Section~\ref{sec:repo} show how much a closure can differ at fixed
    conductivity.
  \item \textbf{Non-identifiability.} Different closures fit the same data
    when the data do not excite the directions in which they differ, such
    as frequencies near the collision rate.
  \item \textbf{Entanglement with numerics.} A correction trained at one
    grid and time step absorbs discretization error, and is then a
    different model at another resolution.
  \item \textbf{Leakage.} Test cases that influenced normalization,
    hyperparameters or early stopping.
\end{itemize}

\subsection{What must be declared}

A comparison between closures is meaningful only at a matched
\emph{information budget}: which equations, paired fields, ensembles,
cumulants, histories, closure targets and solver calls each method
received (see \texttt{model-bakeoff/docs/information\_budgets.md}). Access
to kinetic trajectories or to frequency responses is a different budget
from access to temperature profiles alone, and a method given more will
appear better for that reason.
